% META: {"id": "TSPE-GEOESPACE-ME-009", "chapitre": "TSPE-GEOMETRIE-ESPACE", "type_objet": "methode", "capacites_codes": ["C15"], "methodes": ["M9"], "sources_inspiration": [], "status": "generated"}
% BEGIN-VERIFY
% from sympy import Matrix, symbols, solve, Eq
% t = symbols('t')
% n=Matrix([1,1,1]); u=Matrix([1,-1,0]); P=Matrix([0,0,0])
% assert n.dot(u) == 0                          # direction parallele au plan
% assert n.dot(P) - 3 != 0                      # point hors du plan
% assert solve(Eq(n.dot(P + t*u) - 3, 0), t) == []
% END-VERIFY
\begin{fichemethode}{M9}{Traduire une intersection en système et l'interpréter}

\quandUtiliser{%
  Pour déterminer l'intersection d'une droite et d'un plan, ou de deux plans,
  et surtout pour \emph{interpréter} le résultat algébrique en position
  relative. La conclusion attendue est géométrique, pas seulement numérique.
}

\pasApas{%
  \begin{enumerate}
    \item \textbf{Écrire} la représentation paramétrique de la droite et
          l'équation cartésienne du plan.
    \item \textbf{Substituer} les expressions de $x$, $y$, $z$ dans l'équation
          du plan~: on obtient une équation d'inconnue $t$.
    \item \textbf{Résoudre} sans préjuger du nombre de solutions.
    \item \textbf{Interpréter}~:
          une solution unique, la droite \emph{perce} le plan~;
          aucune solution, elle lui est \emph{strictement parallèle}~;
          une infinité, elle est \emph{contenue} dans le plan.
    \item \textbf{Confirmer} avec les vecteurs~: si
          $\vec{n}\cdot\vec{u}=0$, la droite est parallèle au plan, et c'est
          l'appartenance d'un point qui départage les deux derniers cas.
  \end{enumerate}
}

\verifier{%
  Dans le cas d'une solution unique, reporter le point trouvé dans l'équation
  du plan \emph{et} dans la représentation de la droite~: il doit satisfaire
  les deux. Dans les autres cas, le test $\vec{n}\cdot\vec{u}=0$ fournit une
  confirmation indépendante.
}

\pieges{%
  Un système sans solution n'est pas une erreur de calcul~: c'est la réponse,
  et elle signifie que les objets ne se rencontrent pas. De même, l'équation
  $0=0$ n'est pas une impasse~: elle traduit une infinité de points communs.
}

\exempleRedige{%
  $\mathcal{P}:\ x+y+z-3=0$ et $d:\ x=t,\ y=-t,\ z=0$.
  En substituant~: $t-t+0-3=0$, soit $-3=0$. Le système n'a aucune solution.
  On confirme avec les vecteurs~: $\vec{u}(1;-1;0)$ vérifie
  $\vec{n}\cdot\vec{u}=1-1+0=0$, et le point $(0;0;0)$ de $d$ ne vérifie pas
  l'équation du plan. Donc $d$ est strictement parallèle à $\mathcal{P}$.
}

\end{fichemethode}
